\documentclass[10pt, letterpaper]{article}

% Packages:
\usepackage[
    ignoreheadfoot, % set margins without considering header and footer
    top=2 cm, % seperation between body and page edge from the top
    bottom=2 cm, % seperation between body and page edge from the bottom
    left=2 cm, % seperation between body and page edge from the left
    right=2 cm, % seperation between body and page edge from the right
    footskip=1.0 cm, % seperation between body and footer
]{geometry} % for adjusting page geometry
\usepackage{titlesec} % for customizing section titles
\usepackage{tabularx} % for making tables with fixed width columns
\usepackage{array} % tabularx requires this
\usepackage[dvipsnames]{xcolor} % for coloring text
\definecolor{primaryColor}{RGB}{0, 79, 144} % define primary color
\usepackage{enumitem} % for customizing lists
\usepackage{fontawesome5} % for using icons
\usepackage{amsmath} % for math
\usepackage[
    pdftitle={Naramala Mani Vardhan Naidu's CV - Mechanical Engineering},
    pdfauthor={Naramala Mani Vardhan Naidu},
    pdfcreator={LaTeX with RenderCV},
    colorlinks=true,
    urlcolor=primaryColor
]{hyperref} % for links, metadata and bookmarks
\usepackage[pscoord]{eso-pic} % for floating text on the page
\usepackage{calc} % for calculating lengths
\usepackage{bookmark} % for bookmarks
\usepackage{lastpage} % for getting the total number of pages
\usepackage{changepage} % for one column entries (adjustwidth environment)
\usepackage{paracol} % for two and three column entries
\usepackage{ifthen} % for conditional statements
\usepackage{needspace} % for avoiding page brake right after the section title
\usepackage{iftex} % check if engine is pdflatex, xetex or luatex

% Ensure that generate pdf is machine readable/ATS parsable:
\ifPDFTeX
    \input{glyphtounicode}
    \pdfgentounicode=1
    \usepackage[utf8]{inputenc}
    \usepackage{lmodern}
\fi

% Some settings:
\AtBeginEnvironment{adjustwidth}{\partopsep0pt} % remove space before adjustwidth environment
\pagestyle{empty} % no header or footer
\setcounter{secnumdepth}{0} % no section numbering
\setlength{\parindent}{0pt} % no indentation
\setlength{\topskip}{0pt} % no top skip
\setlength{\columnsep}{0cm} % set column seperation
\makeatletter
\let\ps@customFooterStyle\ps@plain % Copy the plain style to customFooterStyle
\patchcmd{\ps@customFooterStyle}{\thepage}{
    \color{gray}\textit{\small Naramala Mani Vardhan Naidu - Page \thepage{} of \pageref*{LastPage}}
}{}{} % replace number by desired string
\makeatother
\pagestyle{customFooterStyle}

\titleformat{\section}{\needspace{4\baselineskip}\bfseries\large}{}{0pt}{}[\vspace{1pt}\titlerule]

\titlespacing{\section}{
    -1pt
}{
    0.3 cm
}{
    0.2 cm
} % section title spacing

\renewcommand\labelitemi{$\circ$} % custom bullet points
\newenvironment{highlights}{
    \begin{itemize}[
        topsep=0.10 cm,
        parsep=0.10 cm,
        partopsep=0pt,
        itemsep=0pt,
        leftmargin=0.4 cm + 10pt
    ]
}{
    \end{itemize}
} % new environment for highlights

\newenvironment{highlightsforbulletentries}{
    \begin{itemize}[
        topsep=0.10 cm,
        parsep=0.10 cm,
        partopsep=0pt,
        itemsep=0pt,
        leftmargin=10pt
    ]
}{
    \end{itemize}
} % new environment for highlights for bullet entries

\newenvironment{onecolentry}{
    \begin{adjustwidth}{
        0.2 cm + 0.00001 cm
    }{
        0.2 cm + 0.00001 cm
    }
}{
    \end{adjustwidth}
} % new environment for one column entries

\newenvironment{twocolentry}[2][]{
    \onecolentry
    \def\secondColumn{#2}
    \setcolumnwidth{\fill, 4.5 cm}
    \begin{paracol}{2}
}{
    \switchcolumn \raggedleft \secondColumn
    \end{paracol}
    \endonecolentry
} % new environment for two column entries

\newenvironment{header}{
    \setlength{\topsep}{0pt}\par\kern\topsep\centering\linespread{1.5}
}{
    \par\kern\topsep
} % new environment for the header

\newcommand{\placelastupdatedtext}{%
  \AddToShipoutPictureFG*{%
    \put(
        \LenToUnit{\paperwidth-2 cm-0.2 cm+0.05cm},
        \LenToUnit{\paperheight-1.0 cm}
    ){\vtop{{\null}\makebox[0pt][c]{
        \small\color{gray}\textit{Last updated in December 2025}\hspace{\widthof{Last updated in December 2025}}
    }}}
  }%
}%

% save the original href command in a new command:
\let\hrefWithoutArrow\href

% new command for external links:
\renewcommand{\href}[2]{\hrefWithoutArrow{#1}{\ifthenelse{\equal{#2}{}}{ }{#2 }\raisebox{.15ex}{\footnotesize \faExternalLink*}}}

\begin{document}
    \newcommand{\AND}{\unskip
        \cleaders\copy\ANDbox\hskip\wd\ANDbox
        \ignorespaces
    }
    \newsavebox\ANDbox
    \sbox\ANDbox{}

    \placelastupdatedtext
    \begin{header}
        \textbf{\fontsize{24 pt}{24 pt}\selectfont Naramala Mani Vardhan Naidu}

        \vspace{0.3 cm}

        \normalsize
        \mbox{{\color{black}\footnotesize\faPhone*}\hspace*{0.13cm}7995115709}%
        \kern 0.25 cm%
        \AND%
        \kern 0.25 cm%
        \mbox{\hrefWithoutArrow{mailto:manivardhan306@gmail.com}{\color{black}{\footnotesize\faEnvelope[regular]}\hspace*{0.13cm}manivardhan306@gmail.com}}%
        \kern 0.25 cm%
        \AND%
        \kern 0.25 cm%
        \mbox{{\color{black}\footnotesize\faMapMarker*}\hspace*{0.13cm}Visakhapatnam}%
    \end{header}

    \vspace{0.3 cm - 0.3 cm}

    \section{Career Objective}
    \begin{onecolentry}
        Motivated and fast-learning B.Tech graduate seeking an entry-level opportunity in IT / Software / Data roles. Possessing strong foundational knowledge in Python, programming logic, problem-solving, and tools used in engineering and analytics. Eager to contribute to organizational growth while upgrading technical skills and gaining practical industry experience in data analysis and software development.
    \end{onecolentry}

    \section{Education}
    \begin{twocolentry}{
        \textit{Sept 2021 – Apr 2025}
    }
        \textbf{B.Tech in Mechanical Engineering} \\
        \textit{GITAM University, Visakhapatnam} \\
        \textit{CGPA: 7.14/10}
    \end{twocolentry}

    \vspace{0.10 cm}
    \begin{twocolentry}{
        \textit{June 2019 – July 2021}
    }
        \textbf{Higher Secondary Education (12th Grade)} \\
        \textit{Sri Chaitanya Junior College} \\
        \textit{Percentage: 87\%}
    \end{twocolentry}

    \vspace{0.10 cm}
    \begin{twocolentry}{
        \textit{June 2018 – Apr 2019}
    }
        \textbf{Secondary School Certificate (10th Class)} \\
        \textit{Sri Chaitanya Techno School} \\
        \textit{CGPA: 9.0/10}
    \end{twocolentry}

    \section{Internship Experience}
    \begin{twocolentry}{
        \textit{Visakhapatnam} \\
        \textit{May 2024 – June 2024}
    }
        \textbf{Internship at Visakhapatnam Steel Plant} \\
        \textit{Engineering Shops and Foundry Department (ES\&F)}
    \end{twocolentry}
    \vspace{0.10 cm}
    \begin{onecolentry}
        \begin{highlights}
            \item Completed a 4-week internship in the ES\&F Department, gaining hands-on experience in material selection, cutting, welding, TIG/spot welding, and machining for equipment spares.
            \item Assisted in component preparation, preliminary checks, and understanding quality procedures and safety protocols.
            \item Developed strong knowledge of workflow coordination between design, planning, and fabrication teams along with industrial standards and documentation practices.
        \end{highlights}
    \end{onecolentry}

    \section{Projects}
    \begin{twocolentry}{}
        \textbf{Basics in Fabrication of Equipment}
    \end{twocolentry}
    \vspace{0.10 cm}
    \begin{onecolentry}
        \textit{Industrial Project – Visakhapatnam Steel Plant}
        \begin{highlights}
            \item Fabricated and learned manufacturing skills in an industrial environment, gaining practical knowledge in equipment fabrication techniques and processes.
        \end{highlights}
    \end{onecolentry}

    \vspace{0.2 cm}
    \begin{twocolentry}{}
        \textbf{BAJA SAE 2023}
    \end{twocolentry}
    \vspace{0.10 cm}
    \begin{onecolentry}
        \begin{highlights}
            \item Led the Braking Department in designing, developing, and optimizing a reliable and high-performance braking system for the BAJA SAE India ATV.
            \item Oversaw the selection, design, and integration of brake components---including calipers, discs, pedals, and hydraulic lines---using CAD tools to ensure precision and manufacturability.
            \item Coordinated testing, simulations, and on-vehicle validation to enhance braking efficiency, safety, heat management, and overall durability under race conditions.
        \end{highlights}
    \end{onecolentry}

    \vspace{0.2 cm}
    \begin{twocolentry}{}
        \textbf{Compression Behaviour in Aluminium Based Composite Foams}
    \end{twocolentry}
    \vspace{0.10 cm}
    \begin{onecolentry}
        \textit{Academic Project}
        \begin{highlights}
            \item Partnered with professor to study the morphology and compression properties of sintered Al/SiC composite foams (CFs) with varying reinforcement (5\%, 10\%, and 15\%) of SiC particles.
            \item Discovered that SiC particles were irregular and angular in shape while aluminium powder had elliptically shaped primary particles. Compression strength of foams increased with increase in silicon carbide percentage.
            \item Found that the addition of SiC particles influenced the viscosity of the melt, leading to thickening of cell walls and more uniform pore size distribution. SiC particles improved wettability and interface bonding between matrix and reinforcement, resulting in reduced necking phenomenon.
        \end{highlights}
    \end{onecolentry}

    \section{Technical Skills}
    \begin{onecolentry}
        \textbf{Languages:} Python \\
        \textbf{Data \& Analytics:} Basic Excel, Data Handling, Problem Solving \\
        \textbf{IT Skills:} Basic SQL, MS Office, Power BI, Logical Thinking
    \end{onecolentry}

    \section{Certifications}
    \begin{onecolentry}
        \begin{highlights}
            \item Basics in Fabrication and Equipment – Visakhapatnam Steel Plant
            \item C and Python Courses – FacePrep CRT
            \item BAJA SAE 2023 – SAE India
            \item Internship at Centre of Excellence in Maritime and Shipbuilding
        \end{highlights}
    \end{onecolentry}

    \section{Languages}
    \begin{onecolentry}
        English, Telugu
    \end{onecolentry}

\end{document}
